\chapter{はじめに}
\section{はじめに}
本テンプレートは東京工科大学大学院\cite{teu}のバイオ・情報メディア研究科
コンピュータサイエンス専攻の修士論文用に作成されました．以下にテンプレー
トの使い方を簡単に説明していきます．また，{\TeX}全般に関する使い方は，奥
村先生の美文書作成入門\cite{okumura}を参考にしてください．

\section{論文情報の登録}
thesis.tex を開くと，ファイル上部に``\verb|\title{論文タイトル}|''や
``\verb|\author{工~~科~~~太~~郎}|''など論文の情報を入力する箇所があります．
まずは，これらの情報を自身のものに適宜変更してください．本テンプレートで
は入力された情報から``受領書''や``ファイルのカバー''，``内表紙''，``概
要''，``英語要旨''を作成しますので，間違えないように気をつけて入力してく
ださい．

\section{本テンプレートで用意したコマンド}
本テンプレートでは図や表を参照するためのコマンドを用意しました．

\subsection{図の例}
図の参照には\verb|\figref{}|を用います．このコマンドは，図を参照する際に
``図''という文字を自動的に付与します．また，1回目の参照では\textbf{太字}
で表示し，2回目以降は通常表示します．

例）本学園の校章を\figref{fig:teu}に示す．
\begin{figure}[tb]
 \centering
 \includegraphics[clip, width=0.5\linewidth]{teulogo_new.eps}
  \caption{本学園の校章}
 \label{fig:teu}
\end{figure}

\subsection{表の例}
表の参照には\verb|\tabref{}|を用います．このコマンドは，表を参照する際に
``表''という文字を自動的に付与します．また，1回目の参照では\textbf{太字}
で表示し，2回目以降は通常表示します．

例）東京工科大学の学部別の学生数をまとめた表を\tabref{tab:teu}に示す．
\begin{table}[tb]
 \begin{center}
  \caption{学部別の人数（2009年5月1日現在）} \label{tab:teu}
  \begin{tabular}{l|c|c|c}\hline\hline
   学部 & 男子（名） & 女子（名） & 合計（名）\\\hline
   メディア学部 & 1618 & 553 & 2151\\
   応用生物学部 & 1081 & 417 & 1498\\
   コンピュータサイエンス学部 & 2046 & 146 & 2192\\ \hline
  \end{tabular}
 \end{center}
\end{table}

\subsection{式の例}
式の参照には\verb|\eqref{}|を用います．このコマンドは，式を参照する際に
``式''という文字を自動的に付与します．また，1回目の参照では\textbf{太字}
で表示し，2回目以降は通常表示します．

例）オイラーの公式を\eqref{eq:euler}に示す．
\begin{equation}
 e^{i\theta} = \cos\theta + i \sin\theta\label{eq:euler}
\end{equation}

\subsection{プログラムの例}
プログラムの参照には\verb|\pgref{}|を用います．このコマンドは，プログラム
を参照する際に``プログラム''という文字を自動的に付与します．また，1回目の
参照では\textbf{太字} で表示し，2回目以降は通常表示します．

例）プログラムの表示の方法には2種類あります．1つは，外部に保存したファイル
を読み込んで表示する方法です．もう１つは，直接プログラムを記述する方法で
す．

``Hello World''を表示するCのプログラムを\pgref{pg:helloc}に示す．
\lstinputlisting[caption=Hello World, label=pg:helloc,
language=c]{hello.c}

``Hello Java''を表示するJavaのプログラムを\pgref{pg:hellojava}に示す．
\begin{lstlisting}[caption=Hello Java, label=pg:hellojava,language=java]
 public class HelloJava {
     public static void main(String[] args) {
         System.out.println("Hello Java");
     }
 }
\end{lstlisting}
